\subsection{The two norm calculations in the geometric freeness proof}
\label{subsec:faithful-free-norm-calculations}

\begin{srlemma}[The integral coefficient lattice]
\label{lem:faithful-integral-lattice-norm}
Let \(E\) be a number field and \(O=\mathcal O_E\).
Let \(\Lambda\) be an \(O\)-module whose underlying abelian group
is finite free, and put \(s=\operatorname{rank}_O\Lambda\).
For every \(a\in O\), multiplication \(m_a\) on \(\Lambda\),
regarded as an integral endomorphism, satisfies
\[
 \det_{\mathbb Z}(m_a)=N_{E/\mathbb Q}(a)^s,\qquad
 \operatorname{charpol}_{\mathbb Z}(m_a)
                   =\operatorname{charpol}_{\mathbb Z}(a|O)^s .
 \tag{NN1}\label{eq:faithful-integral-lattice-character}
\]
If \(a\ne0\), \(k\) is any field of characteristic the prime \(p\),
and \(m_{a,k}\) is its base change to \(k\otimes_{\mathbb Z}\Lambda\),
then
\[
 \dim_k\ker m_{a,k}\le s\,v_p(|N_{E/\mathbb Q}(a)|).
 \tag{NN2}\label{eq:faithful-integral-lattice-kernel}
\]
The characteristic polynomial of \(m_{a,k}\) is the image of
\eqref{eq:faithful-integral-lattice-character} in \(k[T]\).
In particular \(\dim_k(k\otimes_{\mathbb Z}\Lambda)=[E:\mathbb Q]s\).
\end{srlemma}
\begin{proof}
We follow \srcid{KernelRankNormBound}, without assuming that
\(\Lambda\) is free over \(O\).
Choose \(s\) independent elements over \(O\), giving an injective
\(O\)-linear map \(j:O^s\to\Lambda\).
Its cokernel is \(O\)-torsion.
Every nonzero \(b\in O\) divides a nonzero integer: a polynomial
relation for \(b\) with nonzero constant term gives such a divisor
by rearranging the relation.
Consequently the cokernel is also integral torsion, and \(j\)
is an inclusion of integral lattices of equal rank.
It intertwines multiplication by \(a\).
Choose integral bases.  The matrix of \(j\) has nonzero determinant,
so taking determinants in
\(j\,m_a|O^s=m_a|\Lambda\,j\) and cancelling that determinant
shows that the two multiplication determinants agree.
The determinant on \(O^s\) is \(N(a)^s\).
This proves the first formula of \eqref{eq:faithful-integral-lattice-character}.

Apply that determinant formula to \(t-a\) for every integer \(t\).
The characteristic polynomial evaluated at \(t\) is the determinant
of \(tI-m_a\), so it equals \(N(t-a)^s\).
The same identity on \(O\) gives the second formula by polynomial
equality over the infinite ring \(\mathbb Z\).

For completeness, the corank calculation uses the actual Smith
bases of the image of an injective integral endomorphism \(f\).
There are an integral basis \(t_i\) of the target and an integral
basis \(s_i\) of its image with \(s_i=d_i t_i\).
The map \(f\) is, in these bases, a diagonal matrix
\(\operatorname{diag}(d_i)\) followed by a unimodular matrix.
After base change to \(k\), the latter stays invertible.
The kernel dimension is therefore the number of indices for
which \(p\mid d_i\).  The corresponding power of \(p\) divides
\(\prod_i d_i\), and hence divides \(\det f\).
Thus
\[
       \dim_k\ker(f\otimes k)\le v_p(|\det f|).
\]
For \(a\ne0\), its multiplication determinant is nonzero by the
first formula, and this inequality gives
\eqref{eq:faithful-integral-lattice-kernel}.
Characteristic polynomials commute with this base change.
Their degrees give the final rank equality.
\end{proof}

The next lemma is the arithmetic end of
\srcid{DegreeMultiplicationNorm}.  It is stated explicitly because
the formal proof obtains the norm-degree equality from a coordinate
growth bound; it does not assume that equality as a property of
the \(O\)-action.

\begin{srlemma}[Multiplicative degree and the integral-coordinate bound]
\label{lem:faithful-degree-function-norm}
Let \(O=\mathcal O_E\) have integral basis
\(b_1,\ldots,b_n\), \(n=[E:\mathbb Q]>0\).
Let \(d:O\to\mathbb N\), \(e,C\in\mathbb N\), and assume:
\begin{enumerate}
\item \(d(ab)=d(a)d(b)\) for nonzero \(a,b\), and
\(d(a)\ge1\) for \(a\ne0\);
\item \(d(m)=m^e\) for every positive integer \(m\);
\item whenever \(a=\sum_i c_i b_i\ne0\) with
\(|c_i|\le T\), one has \(d(a)\le C T^e\).
\end{enumerate}
Then for every \(a\ne0\),
\[
                        d(a)^n=|N_{E/\mathbb Q}(a)|^e .
 \tag{NN3}\label{eq:faithful-degree-function-norm}
\]
\end{srlemma}
\begin{proof}
First \(d(a^k)=d(a)^k\), including \(k=0\), by (1) and
\(d(1)=1\).
Write \(M=|N(a)|\ne0\).
For every \(k\ge0\), let \(T\) be the least nonnegative integer
with
\[
                    M^k<(T+1)^n.
\]
Then \(T^n\le M^k\).
The quotient \(O/(a^k)\) has \(M^k\) elements.
The \((T+1)^n\) coordinate vectors
\((u_i)\in\{0,\ldots,T\}^n\) cannot all have distinct images
under \((u_i)\mapsto\sum_i u_i b_i\bmod(a^k)\).
Subtract two distinct vectors with the same image.
Their difference gives
\[
             0\ne x=\sum_i c_i b_i=y a^k,\qquad |c_i|\le T.
\]
Independence of the integral basis makes \(x\ne0\), and hence
\(y\ne0\).
By (1) and (3),
\[
                       d(a)^k\le d(y)d(a)^k=d(x)\le CT^e .
\]
Raise this inequality to the \(n\)-th power and use \(T^n\le M^k\):
\[
                  (d(a)^n)^k\le C^n(M^e)^k .
\]
If \(d(a)^n>M^e\), the ratio of these two positive integer bases
would have unbounded powers, contradicting this fixed constant
for all \(k\).  Therefore
\[
                          d(a)^n\le M^e .
 \tag{NN4}\label{eq:faithful-degree-norm-upper}
\]

The integral ideal \((a)\) contains its absolute norm \(M\).
Thus \(M=ba\) for some nonzero \(b\in O\).
By (1)--(2),
\[
                  d(b)d(a)=M^e.
\]
Taking algebraic norms in \(ba=M\) gives
\[
                    |N(b)|M=M^n.
\]
Consequently
\[
              d(b)^n d(a)^n=|N(b)|^e M^e.
\]
Apply \eqref{eq:faithful-degree-norm-upper} to \(b\) and cancel
the positive factor \(|N(b)|^e\).  This gives
\(M^e\le d(a)^n\), proving
\eqref{eq:faithful-degree-function-norm}.
\end{proof}

\begin{srlemma}[The effective quadratic bound used for degree growth]
\label{lem:faithful-effective-quadratic-bound}
Let \(B,P\) be abelian groups and \(S\subset P\) an additive
submonoid.  Suppose \(q:B\to P\) has \(q(0)=0\) and
\[
                q(x+y)+q(x-y)=2q(x)+2q(y).
\]
Let \(h\in S\), suppose \(q(x)\in S\) for all \(x\), and suppose
that for each \(x\) there is \(N\ge0\) with \(Nh-q(x)\in S\).
For any finite list \(\omega_1,\ldots,\omega_m\in B\) there is
\(C\ge1\) such that
\[
  CT^2h-q\!\left(\sum_{i=1}^m c_i\omega_i\right)\in S
                    \qquad (|c_i|\le T,\ c_i\in\mathbb Z).
 \tag{NN5}\label{eq:faithful-effective-quadratic-bound}
\]
\end{srlemma}
\begin{proof}
The parallelogram identity gives \(q(-x)=q(x)\) and, by
two-step induction, \(q(nx)=n^2q(x)\) for natural \(n\);
hence \(q(cx)=c^2q(x)\) for all integers \(c\).
Induct on the finite list.
For the empty list take \(C=1\).
Suppose the tail has constant \(C\), and choose \(N\) with
\(Nh-q(\omega)\in S\) for the added element.
For \(|c|\le T\), both
\(NT^2h-q(c\omega)\) and the tail expression
\(CT^2h-q(y)\) lie in \(S\): for the first use
\(q(c\omega)=c^2q(\omega)\) and \(T^2-c^2\ge0\).
The parallelogram identity gives
\[
 (2N+2C)T^2h-q(c\omega+y)
  =2\bigl(NT^2h-q(c\omega)\bigr)
   +2\bigl(CT^2h-q(y)\bigr)+q(c\omega-y).
\]
All three terms belong to \(S\).
Use \(2N+2C\) as the new positive constant.
\end{proof}

\paragraph{Frozen-source ledger and geometric interface.}
Lemma~\ref{lem:faithful-integral-lattice-norm} is
\href{https://github.com/ChristianKudeba/alpoge-rh-verification/blob/alpoge-source-7cbf7a0c/src/Development/WeilBound/CoreStatements/TorusUniformisation/ChowMahlerFamilies/KernelRankNormBound.lean}
{\srcid{KernelRankNormBound}}
(blob \srcid{731390fb9ea36fcdfffd4246d81834d09f5b05ab}):
\srcid{det\_eq\_of\_intertwine}, \srcid{det\_lsmul},
\srcid{charpoly\_lsmul}, \srcid{corank\_le}, and the final theorem.
Lemmas~\ref{lem:faithful-degree-function-norm} and
\ref{lem:faithful-effective-quadratic-bound} are respectively
\srcid{pow\_eq\_norm} and \srcid{quad\_bound} (with
\srcid{quad\_nat}, \srcid{quad\_int}, and
\srcid{le\_of\_pow\_le}) in
\href{https://github.com/ChristianKudeba/alpoge-rh-verification/blob/alpoge-source-7cbf7a0c/src/Development/WeilBound/CoreStatements/TorusUniformisation/ChowMahlerFamilies/DegreeMultiplicationNorm.lean}
{\srcid{DegreeMultiplicationNorm}}
(blob \srcid{09a4b78a396ac292ab8fe3938f653e00cca99a6a}).
Both source files were checked against the frozen baseline.

The geometric application takes \(d(a)\) to be the constant
finite-flat rank of \([a]\) on an abelian variety of dimension
\(g\), with \(e=2g\).
Its remaining source steps establish finiteness, multiplicativity,
the integer value \(d(m)=m^{2g}\), and the coordinate growth
\(CT^{2g}\) by symmetric very ample bundles and section-rank
sandwiches.  Those steps are not assumed proved by the arithmetic
lemmas above and are recorded separately in the geometric degree
producer.
